\documentclass[11pt, a4paper]{article}
\usepackage[utf8]{inputenc}
\usepackage[margin=1in]{geometry}
\usepackage{longtable}
\usepackage{makecell}
\usepackage{array}
\usepackage{graphicx}

% Formatting lengths
\setlength{\parindent}{0pt}
\setlength{\parskip}{0.8em}
\renewcommand{\arraystretch}{1.5}

\begin{document}

% Header Section
\noindent
\begin{minipage}[t]{0.48\textwidth}
\vspace{0pt}
\includegraphics[width=5cm]{smartbridge_logo.png}
\end{minipage}
\hfill
\begin{minipage}[t]{0.48\textwidth}
\vspace{0pt}
\raggedleft
\includegraphics[width=4.2cm]{skillwallet_logo.png}
\end{minipage}

% Project Details Table
\begin{center}
\begin{tabular}{|p{0.30\textwidth}|p{0.60\textwidth}|}
\hline
\textbf{Date} & 15 March 2026 \\ \hline
\textbf{Team ID} & XXXXXX \\ \hline
\textbf{Project Name} & Rising Waters: A Machine Learning Approach to Flood Prediction \\ \hline
\textbf{Maximum Marks} & 5 Marks \\ \hline
\end{tabular}
\end{center}

\vspace{1em}

% Title and Intro
\begin{center}
    \Large \textbf{Coding \& Solution}
\end{center}

This section evaluates the quality and completeness of your implemented code and the overall solution. Your submission should demonstrate working functionality, clean code practices, and alignment with your defined requirements.

\textbf{Instructions:}
\begin{enumerate}
    \item Ensure your code is well-structured, commented, and follows best practices for your chosen technology stack.
    \item The solution should address the core problem statement and implement the key functional requirements.
    \item Include all source files, configuration files, and setup instructions needed to run the application.
    \item Attach or link to your code repository (e.g., GitHub) in the table below.
\end{enumerate}

\vspace{1em}

% Solution Summary Table
\textbf{Solution Summary}

\noindent
\begin{longtable}{|>{\bfseries}p{0.30\textwidth}|p{0.65\textwidth}|}
\hline
Field & \textbf{Details} \\ \hline
\endfirsthead
\endhead % Empty to prevent repeating headers on new pages
\endfoot % Empty to prevent repeating footers on new pages
\endlastfoot

Repository Link / URL & [Paste your GitHub / GitLab or hosted link here] \\ \hline
Programming Language(s) & Python, HTML5, CSS3, JavaScript \\ \hline
Framework(s) Used & Flask, Bootstrap 5, XGBoost, Scikit-learn \\ \hline
Key Features Implemented & Flood prediction, XGBoost model, Flask web app, Risk analysis, Safety recommendations, Responsive UI \\ \hline
Pending / Incomplete Features & Live weather API integration, Interactive map visualization (Future Scope) \\ \hline
Setup / Run Instructions & Install requirements → Run python app.py → Open http://127.0.0.1:5000 \\ \hline
\end{longtable}

\vspace{1em}

% Code Quality Checklist Table
\textbf{Code Quality Checklist}

\noindent
\begin{longtable}{|c|>{\raggedright\arraybackslash}p{0.65\textwidth}|c|}
\hline
\textbf{S.No} & \textbf{Criteria} & \textbf{Status (Yes / No)} \\ \hline
\endfirsthead
\endhead
\endfoot
\endlastfoot

1 & Code is modular and organized into functions / classes & Yes \\ \hline
2 & Meaningful variable and function names are used & Yes \\ \hline
3 & Code includes comments / documentation where necessary & Partial \\ \hline
4 & Error handling is implemented for critical operations & Yes \\ \hline
5 & The application runs without critical errors & Yes \\ \hline
6 & Code is committed to a version control repository & Yes \\ \hline
\end{longtable}

\vspace{1em}

\textbf{Additional Notes / Comments:}

The project successfully implements an AI-powered flood prediction system using the XGBoost Machine Learning algorithm. The application follows a modular architecture with a Flask backend and responsive web interface, providing accurate flood probability prediction and risk analysis.

\end{document}